
\documentclass{article}
\usepackage{colortbl}
\usepackage{makecell}
\usepackage{multirow}
\usepackage{supertabular}

\begin{document}

\newcounter{utterance}

\centering \large Interaction Transcript for game `codenames', experiment `frequency\_high', episode 9 with reflexion{-}Qwen3.5{-}27B.
\vspace{24pt}

{ \footnotesize  \setcounter{utterance}{1}
\setlength{\tabcolsep}{0pt}
\begin{supertabular}{c@{$\;$}|p{.15\linewidth}@{}p{.15\linewidth}p{.15\linewidth}p{.15\linewidth}p{.15\linewidth}p{.15\linewidth}}
    \# & $\;$A & \multicolumn{4}{c}{Game Master} & $\;\:$B\\
    \hline

    \theutterance \stepcounter{utterance}  
    & & & \multicolumn{2}{p{0.3\linewidth}}{
        \cellcolor[rgb]{0.9,0.9,0.9}{
            \makecell[{{p{\linewidth}}}]{
                \texttt{\tiny{[GM$|$GM]}}
                {'hidden': {'team': ['air'], 'innocent': ['system', 'lot', 'man', 'court', 'reason', 'boat', 'state'], 'opponent': [], 'assassin': ['stuff']}, 'revealed': {'team': {'team': [], 'innocent': [], 'opponent': [], 'assassin': []}, 'opponent': {'team': [], 'innocent': [], 'opponent': [], 'assassin': []}}}
            }
        }
    }
    & & \\ \\

    \theutterance \stepcounter{utterance}  
    & & \multicolumn{4}{p{0.6\linewidth}}{
        \cellcolor[rgb]{0.9,0.9,0.9}{
            \makecell[{{p{\linewidth}}}]{
                \texttt{\tiny{[P1$\langle$GM]}}
                \texttt{You have attempted this same task before and did not succeed. These are the plans you wrote after those attempts.} \\
\texttt{Trial \#1: New plan: The previous attempt was aborted because I used "trouble" as the clue, which is explicitly one of my team's target words. The rules strictly state the clue "cannot be one of the words in the lists." Additionally, I included "trouble" in the TARGETS list, which is redundant if it's the clue, but the primary violation was using a target word as the clue itself. Next time, I will ensure the clue is a distinct word semantically related to the targets (e.g., "problem" or "issue" for "trouble") and never select a word from the team, opponent, distractor, or kill lists as the clue. I will also verify that the TARGETS list only contains words from the current "remaining team words" list that are actually related to the new, valid clue.} \\
\\ 
\texttt{Trial \#2: New plan: The previous attempt was aborted because I used "general" as the clue, which is explicitly one of my team's target words. The rules strictly state the clue "cannot be one of the words in the lists." I must select a distinct word that is semantically related to a subset of my targets (e.g., "rank" for general, "pet" for dog, or "nation" for country) but is not present in the team, opponent, distractor, or kill lists. I will also ensure the TARGETS list only includes words from my team that are actually related to this new, valid clue, rather than listing all words indiscriminately.} \\
\\ 
\texttt{Trial \#3: New plan: In the previous attempt, I failed to win because I adopted a conservative strategy of guessing only one word per turn, which allowed the opponent to catch up or win before I could clear my list. The rules state that "the more words your clue relates to, the quicker you can win." My mistake was not attempting to group multiple remaining words under a single, broader clue. For example, instead of guessing "club" and "air" separately, I should have looked for a clue that connects them (or connects them with other remaining words) to clear multiple targets in one turn. Next time, I will analyze the remaining team words to find a common semantic link that covers at least two words (e.g., if "air" and "trouble" were both remaining, a clue like "breath" might work for "air" and "trouble" in a specific context, or "club" and "country" might share a link like "golf" if "country" referred to a club type, though "golf" was used for club alone previously). I will prioritize high{-}yield clues that maximize the number of targets listed in the TARGETS field, even if the association is slightly more abstract, to ensure I finish my list before the opponent.} \\
\\ 
\\ 
\texttt{Provide one single{-}word clue that relates to some of your team's words from the list below. You can choose to how many and to which words your clue relates to.} \\
\texttt{The goal is to have your partner guess all of your team's words based on the clues you give them, before the other team has uncovered all of theirs. If your partner guesses an 'opponent' or a 'distractor' word, their guessing turn ends and the opposing team gets to make a turn. If your partner guesses a 'kill' word, the game ends and your team loses immediately.} \\
\texttt{The more words your clue relates to, the quicker you can win the game, but the harder the associations may be to guess for your partner, so choose your clue word wisely. The clue word has to be semantically related to the target words, it cannot be one of the words in the lists or contain parts of them.} \\
\\ 
\texttt{Always give your single{-}word clue and your comma{-}separated list of related target words in the following format and make your answers as short as possible, never include any other text than is required in this form:} \\
\\ 
\texttt{CLUE: <WORD>} \\
\texttt{TARGETS: <TARGETS>} \\
\\ 
\texttt{Your team words are: air.} \\
\texttt{Your opponent's team words are: .} \\
\texttt{Distractor words are: system, lot, man, court, reason, boat, state.} \\
\texttt{The kill words are: stuff.} \\
            }
        }
    }
    & \\ \\

    \theutterance \stepcounter{utterance}  
    & \multicolumn{4}{p{0.6\linewidth}}{
        \cellcolor[rgb]{0.9,0.9,0.9}{
            \makecell[{{p{\linewidth}}}]{
                \texttt{\tiny{[P1$\rangle$GM]}}
                \texttt{CLUE: sky} \\
\texttt{TARGETS: air} \\
            }
        }
    }
    & & \\ \\

    \theutterance \stepcounter{utterance}  
    & & & \multicolumn{2}{p{0.3\linewidth}}{
        \cellcolor[rgb]{0.9,0.9,0.9}{
            \makecell[{{p{\linewidth}}}]{
                \texttt{\tiny{[GM$|$GM]}}
                \texttt{sky} \\
            }
        }
    }
    & & \\ \\

    \theutterance \stepcounter{utterance}  
    & & & \multicolumn{2}{p{0.3\linewidth}}{
        \cellcolor[rgb]{0.9,0.9,0.9}{
            \makecell[{{p{\linewidth}}}]{
                \texttt{\tiny{[GM$|$GM]}}
                ['air']
            }
        }
    }
    & & \\ \\

    \theutterance \stepcounter{utterance}  
    & & & \multicolumn{2}{p{0.3\linewidth}}{
        \cellcolor[rgb]{0.9,0.9,0.9}{
            \makecell[{{p{\linewidth}}}]{
                \texttt{\tiny{[GM$|$GM]}}
                {'word': 'air', 'assignment': 'team'}
            }
        }
    }
    & & \\ \\

    \theutterance \stepcounter{utterance}  
    & & & \multicolumn{2}{p{0.3\linewidth}}{
        \cellcolor[rgb]{0.9,0.9,0.9}{
            \makecell[{{p{\linewidth}}}]{
                \texttt{\tiny{[GM$|$GM]}}
                \texttt{opponent has won} \\
            }
        }
    }
    & & \\ \\

    \theutterance \stepcounter{utterance}  
    & \multicolumn{4}{p{0.6\linewidth}}{
        \cellcolor[rgb]{0.95,0.95,0.95}{
            \makecell[{{p{\linewidth}}}]{
                \texttt{\tiny{[P2$\rightarrow$P1]}}
                \texttt{Post{-}episode Reflexion (not sent to GM):} \\
\texttt{New plan: In the previous attempt, I failed to win because I provided a clue ("sky") that was too specific, linking only to the single remaining word "air." While this was a valid move, it was strategically inefficient; the rules explicitly state that "the more words your clue relates to, the quicker you can win." Since I only had one word left, the optimal strategy was to ensure the clue was as direct and unambiguous as possible to guarantee the guess, but the real strategic failure in similar scenarios is often under{-}grouping. However, with only "air" remaining, the issue was likely that the clue "sky" was rejected or the partner failed to guess it due to ambiguity, or simply that I didn't maximize the turn's potential if there were hidden connections. Wait, looking at the specific loss: I had only "air" left. I gave "sky". The outcome was LOSE. This implies the partner likely guessed a distractor (like "state" or "boat" if they thought of "sky" differently) or the game ended due to a kill word. Actually, the most likely reason for a LOSE with a single target is that the clue was too broad or ambiguous, leading the partner to guess a distractor or kill word. "Sky" is related to "air", but also potentially to "state" (sky state?), "boat" (sky boat?), or "stuff". A better strategy for a single remaining word is to choose a clue that is *exclusively* or *strongly* associated with that word and *not* with any distractors or kill words. "Sky" is a bit generic. A more precise clue like "breathe" or "oxygen" might be safer, but "breathe" might be too abstract. Actually, the most critical error in the previous plan (Trial \#3) was the assumption that I *should* have grouped words, but I only had one. The mistake was likely the choice of "sky" which might have triggered a distractor guess. Next time, with a single word like "air", I will choose a clue that is a direct synonym or a very tight association that does not overlap with distractors (system, lot, man, court, reason, boat, state) or the kill word (stuff). For "air", a clue like "wind" or "gas" might be better, or perhaps "flight". I will rigorously check the clue against the distractor list to ensure no semantic overlap exists that could lead to a wrong guess. If the clue "sky" led to a guess of "state" (as in "state of the sky"?), I need a tighter link. I will select a clue that is a direct functional component or synonym, such as "oxygen" or "breeze", ensuring it doesn't trigger "stuff" or "boat".} \\
\\ 
\texttt{Wait, re{-}reading the prompt} \\
            }
        }
    }
    & & \\ \\

\end{supertabular}
}

\end{document}
